% GENERATED FILE. Edit canonical lessons, then run npm run generate:books.
% canonical-source-hash: fnv1a64:18b8100442455001
% canonical-lessons: TA-C42-come, TA-C42-sit, TA-C42-read, TA-C42-give

\chapter{More Everyday Verbs}
\label{ch:ta-more-verbs}
\hlchaptermodality{42}

\begin{chapteropening}
\textbf{By the end of this chapter:} \emph{I can tell someone to come, sit, read and give in Tamil.}

Complete TA-C42-give and say all four verbs in order.
\end{chapteropening}

\section[\texorpdfstring{\ta{வா} (vā)}{வா (vā)}]{\ta{வா} (vā) — come}

\label{lesson:TA-C42-come}



\begin{quote}
A short run of everyday verbs, one at a time.
\end{quote}

\subsection*{You'll want to know: \ta{வா}}
\textbf{\ta{வா}} (\emph{vā}) — \textquotedblleft{}come\textquotedblright{}.

One syllable, and a complete sentence. The first thing said at a doorway, and the shortest verb form Tamil has: the bare stem, with nothing added.

\textbf{In the family, and next door}

\noindent
\begin{tabularx}{\linewidth}{@{}>{\raggedright\arraybackslash}X>{\raggedright\arraybackslash}X>{\raggedright\arraybackslash}X@{}}
\toprule
\textbf{} & \textbf{word} & \textbf{same root} \\
\midrule
Telugu & \te{రా} (\emph{rā}) &  \\
Malayalam & \ml{വരുക} (\emph{varuka}) & yes \\
Hindi (neighbour) & \dv{आना} (\emph{ānā}) &  \\
English & come &  \\
\bottomrule
\end{tabularx}

\begin{grammarlens}[title={What you've built}]
One more everyday action. Three follow, and each reuses the ones before.
\end{grammarlens}

\subsection*{Your turn}
Say these aloud:

\begin{itemize}
\raggedright
  \item \emph{vā}
  \item it once more, slowly
\end{itemize}

\begin{tcolorbox}[breakable,colback=teal!4,colframe=teal!35!black,title={Before you move on}]
What does \ta{வா} mean? (\textquotedblleft{}come\textquotedblright{}.) Say it once more.
\end{tcolorbox}
\section[\texorpdfstring{\ta{உட்கார்} (uṭkār)}{உட்கார் (uṭkār)}]{\ta{உட்கார்} (uṭkār) — sit}

\label{lesson:TA-C42-sit}



\begin{quote}
Before the new one: what did \emph{vā} mean?
\end{quote}

\subsection*{You'll want to know: \ta{உட்கார்}}
\textbf{\ta{உட்கார்}} (\emph{uṭkār}) — \textquotedblleft{}sit\textquotedblright{}.

Built from \emph{uḷ}, 'inside', and \emph{kār} — to settle oneself in. Tamil often builds a verb from a place-word plus an action this way.

\textbf{In the family, and next door}

\noindent
\begin{tabularx}{\linewidth}{@{}>{\raggedright\arraybackslash}X>{\raggedright\arraybackslash}X@{}}
\toprule
\textbf{} & \textbf{word} \\
\midrule
Kannada & \ka{ಕುಳಿತುಕೋ} (\emph{kuḷitukō}) \\
Telugu & \te{కూర్చో} (\emph{kūrcō}) \\
Hindi (neighbour) & \dv{बैठना} (\emph{baiṭhnā}) \\
English & sit \\
\bottomrule
\end{tabularx}

\begin{grammarlens}[title={What you've built}]
One more everyday action. Three follow, and each reuses the ones before.
\end{grammarlens}

\subsection*{Your turn}
Say these aloud:

\begin{itemize}
\raggedright
  \item \emph{uṭkār}
  \item it once more, slowly
  \item it after \emph{vā}, so the two sit together
\end{itemize}

\begin{tcolorbox}[breakable,colback=teal!4,colframe=teal!35!black,title={Before you move on}]
What does \ta{உட்கார்} mean? (\textquotedblleft{}sit\textquotedblright{}.) Say it once more.
\end{tcolorbox}
\section[\texorpdfstring{\ta{படி} (paṭi)}{படி (paṭi)}]{\ta{படி} (paṭi) — read}

\label{lesson:TA-C42-read}



\begin{quote}
Before the new one: what did \emph{uṭkār} mean?
\end{quote}

\subsection*{You'll want to know: \ta{படி}}
\textbf{\ta{படி}} (\emph{paṭi}) — \textquotedblleft{}read\textquotedblright{}.

Also 'to study'. The same word covers reading a page and studying a subject, so a student \emph{paṭikkiṟār} whether they are holding a book or not.

\textbf{In the family, and next door}

\noindent
\begin{tabularx}{\linewidth}{@{}>{\raggedright\arraybackslash}X>{\raggedright\arraybackslash}X@{}}
\toprule
\textbf{} & \textbf{word} \\
\midrule
Kannada & \ka{ಓದು} (\emph{ōdu}) \\
Telugu & \te{చదువు} (\emph{caduvu}) \\
Malayalam & \ml{വായിക്കുക} (\emph{vāyikkuka}) \\
Hindi (neighbour) & \dv{पढ़ना} (\emph{paṛhnā}) \\
English & read \\
\bottomrule
\end{tabularx}

\begin{grammarlens}[title={What you've built}]
One more everyday action. Three follow, and each reuses the ones before.
\end{grammarlens}

\subsection*{Your turn}
Say these aloud:

\begin{itemize}
\raggedright
  \item \emph{paṭi}
  \item it once more, slowly
  \item it after \emph{uṭkār}, so the two sit together
\end{itemize}

\begin{tcolorbox}[breakable,colback=teal!4,colframe=teal!35!black,title={Before you move on}]
What does \ta{படி} mean? (\textquotedblleft{}read\textquotedblright{}.) Say it once more.
\end{tcolorbox}
\section[\texorpdfstring{\ta{கொடு} (koṭu)}{கொடு (koṭu)}]{\ta{கொடு} (koṭu) — give}

\label{lesson:TA-C42-give}



\begin{quote}
Before the new one: what did \emph{paṭi} mean?
\end{quote}

\subsection*{You'll want to know: \ta{கொடு}}
\textbf{\ta{கொடு}} (\emph{koṭu}) — \textquotedblleft{}give\textquotedblright{}.

The everyday 'give'. Tamil distinguishes giving away from giving toward the speaker; this is the giving-away one.

\textbf{In the family, and next door}

\noindent
\begin{tabularx}{\linewidth}{@{}>{\raggedright\arraybackslash}X>{\raggedright\arraybackslash}X>{\raggedright\arraybackslash}X@{}}
\toprule
\textbf{} & \textbf{word} & \textbf{same root} \\
\midrule
Kannada & \ka{ಕೊಡು} (\emph{koḍu}) & yes \\
Telugu & \te{ఇవ్వు} (\emph{ivvu}) &  \\
Malayalam & \ml{തരൂ} (\emph{tarū}) &  \\
Hindi (neighbour) & \dv{देना} (\emph{denā}) &  \\
English & give &  \\
\bottomrule
\end{tabularx}

\begin{grammarlens}[title={What you've built}]
That closes the run: four more actions, each met on its own.
\end{grammarlens}

\subsection*{Your turn}
Say these aloud:

\begin{itemize}
\raggedright
  \item \emph{koṭu}
  \item it once more, slowly
  \item it after \emph{paṭi}, so the two sit together
\end{itemize}

\begin{tcolorbox}[breakable,colback=teal!4,colframe=teal!35!black,title={Before you move on}]
What does \ta{கொடு} mean? (\textquotedblleft{}give\textquotedblright{}.) Say it once more.
\end{tcolorbox}
